%% ============================================================
%% PREAMBLE / INDSTILLINGER TIL RAPPORT
%% ============================================================

%% --- Encoding og font ---
\usepackage[utf8]{inputenc} % Indtastning af specialtegn, f.eks. æ, ø og å
\usepackage[T1]{fontenc}    % Korrekt visning af specialtegn i PDF-dokumentet

%% --- Skrifttype ---
\usepackage{lmodern}        % Latin Modern (sans-serif)
\renewcommand{\familydefault}{\sfdefault}

%% --- Sidelayout ---
\usepackage{geometry}
\geometry{a4paper, margin=2.5cm}

%% --- Sprog (valgfrit) ---
% Uncomment linjen herunder for danske overskrifter automatisk
% (fx "Figur" i stedet for "Figure", "Tabel" i stedet for "Table").
%  \usepackage[danish]{babel}

%% --- Tekst og formatering ---
\usepackage{lipsum}         % Filler-tekst (kan fjernes når rapporten er færdig)
\usepackage{soul}           % Strikethrough (\st{tekst})
\usepackage{xcolor}         % Farvet tekst (\textcolor{red}{tekst})

%% --- Matematik ---
\usepackage{amsmath}        % Komplekse matematiske ligninger
\usepackage{amssymb}        % Ekstra symboler (checkmark, osv.)

%% --- Tabeller ---
\usepackage{array}          % Ekstra kolonnetyper i tabeller
\usepackage{makecell}       % Line breaks i tabelceller (\makecell{linje1\\linje2})
\usepackage{booktabs}       % Professionelle tabeller (\toprule, \midrule, \bottomrule)
\usepackage{multirow}       % Celler der spænder over flere rækker

%% --- Figurer og billeder ---
\usepackage{graphicx}       % Inkludering af billeder (\includegraphics)
\usepackage{float}          % Præcis placering af figurer ([H])
\usepackage{subcaption}     % Subfigurer side om side (\begin{subfigure})

%% --- Enheder og tal ---
\usepackage{siunitx}        % SI-enheder og talformatering (\SI{10}{\mega\hertz})

%% --- Kode ---
% Minted bruger Pygments til syntax highlighting.
% VIGTIGT: Overleaf kræver at kodefiler har en tekstbaseret filendelse.
% Omdøb .vhd filer til .vhdl for at \inputminted virker korrekt.
\usepackage{minted}
\setminted{
    fontfamily=ttfamily,
    fontsize=\small,
    frame=lines,
    linenos=true,
    tabsize=4,
    obeytabs=true,
    breaklines=true
}

%% --- Hyperlinks ---
\usepackage{hyperref}
\hypersetup{
    colorlinks=true,
    linkcolor=blue!70!black,
    citecolor=black,
    filecolor=black,
    urlcolor=blue!70!black
}

%% --- Kapitel- og sektionsformatering ---
\usepackage{titlesec}
\titleformat{\chapter}[hang]{\normalfont\huge\bfseries}{\thechapter\ }{10pt}{}
\titlespacing*{\chapter}{0pt}{-30pt}{20pt}
% Sektioner uden nummerering i selve teksten (men med numre i indholdsfortegnelsen)
\titleformat{\section}{\normalfont\Large\bfseries}{}{0em}{}
\titleformat{\subsection}{\normalfont\large\bfseries}{}{0em}{}

%% --- Indholdsfortegnelse ---
\renewcommand{\contentsname}{Indhold}

%% --- Header/Footer (valgfrit) ---
% Uncomment for header/footer med sidetal og kapitelnavne:
 \usepackage{fancyhdr}
 \pagestyle{fancy}
 \fancyhf{}
 \fancyhead[L]{\leftmark}
 \fancyhead[R]{62711 -- Project Work F}
 \fancyfoot[C]{\thepage}
 \renewcommand{\headrulewidth}{0.4pt}

% Tekst rundt om billeder/figurer
\usepackage{wrapfig}

\usepackage{circuitikz}